\begin{figure}[htbp]
\centering

\resizebox{0.98\textwidth}{!}{%
\begin{tikzpicture}[
    node distance=5mm and 12mm,
    font=\small,
    data/.style={
        draw,
        rounded corners=2pt,
        align=center,
        text width=4.4cm,
        minimum height=10mm,
        fill=gray!12
    },
    process/.style={
        draw,
        rounded corners=2pt,
        align=center,
        text width=4.4cm,
        minimum height=10mm,
        fill=blue!8
    },
    compute/.style={
        draw,
        rounded corners=2pt,
        align=center,
        text width=4.4cm,
        minimum height=10mm,
        fill=green!10
    },
    optional/.style={
        draw,
        dashed,
        rounded corners=2pt,
        align=center,
        text width=4.4cm,
        minimum height=10mm,
        fill=orange!12
    },
    presentation/.style={
        draw,
        rounded corners=2pt,
        align=center,
        text width=4.4cm,
        minimum height=10mm,
        fill=violet!9
    },
    note/.style={
        draw,
        dashed,
        rounded corners=2pt,
        align=center,
        text width=9.8cm,
        minimum height=8mm,
        fill=yellow!10,
        font=\footnotesize
    },
    flow/.style={
        -{Latex[length=2.2mm]},
        thick
    },
    optionalflow/.style={
        -{Latex[length=2.2mm]},
        thick,
        dashed
    }
]

% Źródło i przygotowanie wejścia
\node[optional] (device) {
    Garmin Venu 3 i Garmin Connect\\
    pomocnicze pozyskanie danych
};

\node[data, below=of device] (workbook) {
    Lokalny workbook \texttt{garmin\_export\_*.xlsx}\\
    formalne wejście systemu
};

\node[process, below=of workbook] (read) {
    Lokalizacja pliku i odczyt XLSX\\
    bez wykonywania obliczeń PWS
};

\node[process, below=of read] (normalize) {
    Ograniczenie struktury i normalizacja\\
    rekordów dziennych oraz aktywności
};

\node[process, below=of normalize] (features) {
    Dzienna oś czasu, cechy dobowe\\
    i agregaty aktywności
};

% Obliczenia pośrednie
\node[compute, below left=7mm and 6mm of features] (load) {
    Metryki obciążenia\\
    okna 7/28 dni, ACWR, EWMA,\\
    monotonia i strain
};

\node[compute, below right=7mm and 6mm of features] (readiness) {
    Gotowość krótkoterminowa\\
    HRV, stres, RHR\\
    i bilans obciążenia
};

\coordinate (componentAnchor)
    at ($(load.south)!0.5!(readiness.south)$);

\node[compute, below=8mm of componentAnchor] (components) {
    Składowe PWS\\
    LQ -- jakość obciążenia,\quad
    REC -- regeneracja,\\
    CAP -- zdolność do utrzymania aktywności,\quad
    VS -- stabilność sygnałów
};

\node[compute, below=of components] (weights) {
    Deterministyczne wyznaczenie wag\\
    kontrola dostępności, wiarygodności\\
    i nakładania informacji
};

\node[compute, below=of weights] (deterministic) {
    Rekord deterministyczny\\
    składowe, wagi, pokrycie i PWS\\
    \textbf{wynik główny systemu}
};

% Opcjonalna ścieżka modeli językowych
\node[optional, right=16mm of deterministic] (llm) {
    Lokalny model językowy lub OpenAI\\
    propozycja zmian wag LQ, REC, CAP i VS
};

\node[optional, below=of llm] (guardrails) {
    Lokalna walidacja i normalizacja\\
    propozycji modelu oraz obliczenie\\
    oddzielnego wariantu PWS
};

% Publikacja i prezentacja
\node[presentation, below=11mm of deterministic] (payload) {
    Payload dashboardu\\
    agregacja gotowych wyników\\
    bez ponownego przeliczania
};

\node[presentation, below=of payload] (snapshot) {
    Publikacja lokalnej migawki\\
    historia, \texttt{current}, walidacja i SHA-256
};

\node[presentation, below=of snapshot] (api) {
    Lokalne API odczytu migawki\\
    \texttt{127.0.0.1:8765/api/v1}
};

\node[presentation, below=of api] (ui) {
    Dashboard webowy\\
    stan bieżący, trendy, pokrycie\\
    i porównanie wariantów PWS
};

\node[note, below=7mm of ui] (audit) {
    W całym przebiegu zachowywane są wersje polityk,
    wartości \texttt{null}, raporty jakości, manifesty,
    błędy oraz informacje o pochodzeniu danych.
};

% Główny przepływ danych
\draw[optionalflow]
    (device) -- node[right, font=\footnotesize]
    {opcjonalne przygotowanie workbooka} (workbook);

\draw[flow] (workbook) -- (read);
\draw[flow] (read) -- (normalize);
\draw[flow] (normalize) -- (features);

\draw[flow] (features) -- (load);
\draw[flow] (features) -- (readiness);

% Gotowość wykorzystuje wcześniej policzony bilans obciążenia
\draw[flow]
    (load.east) -- node[above, font=\footnotesize]
    {bilans obciążenia} (readiness.west);

% Źródła poszczególnych składowych PWS
\draw[flow] (features.south) -- (components.north);
\draw[flow] (load.south) |- (components.west);
\draw[flow] (readiness.south) |- (components.east);

\draw[flow] (components) -- (weights);
\draw[flow] (weights) -- (deterministic);

% Model otrzymuje gotowy kontekst, a nie surowe dane Garmin
\draw[optionalflow]
    (deterministic) --
    node[above, font=\footnotesize]
    {ustrukturyzowany kontekst} (llm);

\draw[optionalflow] (llm) -- (guardrails);

% Wynik główny i warianty pozostają rozdzielone
\draw[flow] (deterministic) -- (payload);
\draw[optionalflow] (guardrails.west) -- (payload.east);

\draw[flow] (payload) -- (snapshot);
\draw[flow] (snapshot) -- (api);
\draw[flow] (api) -- (ui);
\draw[flow] (ui) -- (audit);

\end{tikzpicture}%
}

\caption{Logiczny przepływ danych od lokalnego źródła
do prezentacji wyników}
\label{fig:logical-data-flow}

\figsource{opracowanie własne na podstawie implementacji systemu}
\end{figure}